% ============================================================================
%  Section V — Methodology  (IEEE TIM, IEEEtran)
%  Passive voice / third person throughout.
%  Requires: booktabs, amsmath, tikz + \usetikzlibrary{arrows.meta,positioning}.
% ============================================================================

\section{Methodology}
\label{sec:method}

The processing chain evaluated in this work is deliberately lightweight, so that
it is representative of what an inexpensive sensor node can execute: fixed-length
windowing, a compact set of time-domain features, and a conventional classifier.
The chain is illustrated in Fig.~\ref{fig:pipeline}. The emphasis of this section,
and the principal methodological contribution of the paper, is the set of
\emph{evaluation protocols} used to obtain deployment-realistic performance
estimates.

% ---------------------------------------------------------------------------
\begin{figure*}[t]
\centering
\begin{tikzpicture}[
  font=\footnotesize, >=Latex, node distance=0.9cm,
  box/.style={draw,rounded corners,align=center,minimum height=1.05cm,
              text width=2.5cm,fill=gray!8}]
  \node[box] (sens) {Smartphone MEMS node\\ $100$~Hz, $6$ channels};
  \node[box,right=of sens] (win) {Windowing\\ $2$~s, $50\%$ overlap};
  \node[box,right=of win] (feat) {Time-domain features\\ $14$ per channel};
  \node[box,right=of feat] (clf) {Classifier\\ RF / KNN / SVM};
  \node[box,right=of clf] (out) {Health state\\ (one of six)};
  \draw[->] (sens)--(win); \draw[->] (win)--(feat);
  \draw[->] (feat)--(clf); \draw[->] (clf)--(out);
  \node[draw,dashed,rounded corners,below=0.7cm of feat,text width=9.0cm,
        align=center,fill=blue!4] (eval)
     {\textbf{Evaluation protocols:}\quad RANDOM $\;\mid\;$ GROUP $\;\mid\;$
      CROSS-SPEED $\;\mid\;$ CROSS-LOAD};
  \draw[->,dashed] (clf.south) |- (eval.east);
  \draw[->,dashed] (win.south) |- (eval.west);
\end{tikzpicture}
\caption{Processing chain and evaluation protocols. Identical features are
assessed under four data-partitioning protocols; the protocol, not the model,
is the primary object of study.}
\label{fig:pipeline}
\end{figure*}

\subsection{Segmentation and Windowing}
Let the dataset comprise $R=36$ continuous recordings, each associated with a
health label and an operating condition (supply frequency and load). Every
recording was divided into fixed-length windows of $2$~s with $50\%$ overlap,
matching the segmentation of the source data article. At $100$~Hz this yields
windows of $200$~samples and a total of $32{,}328$ windows, distributed equally
across the six classes ($5{,}388$ per class). The window length and overlap were
treated as configurable parameters and were varied in the trade-off study of
Section~\ref{sec:results}.

\subsection{Feature Extraction}
Each window was reduced to a set of $14$ time-domain features per channel,
summarised in Table~\ref{tab:features}. These comprise standard statistical and
shape descriptors widely used in vibration-based diagnosis, and require no
spectral transform, which keeps the extraction cost compatible with a
microcontroller. For a window of $C$ channels, the per-channel feature vectors
were concatenated into a single descriptor of length $14C$ (for example, $84$ for
the full six-channel set). The six-feature subset marked in
Table~\ref{tab:features}, corresponding to the descriptors used by the source
data article's baseline, was retained for a feature-set ablation. For the
distance- and margin-based classifiers, features were standardised to zero mean
and unit variance using statistics estimated on the training partition only.

\subsection{Classifiers}
Three baseline classifiers were considered, chosen for their prevalence in the
fault-diagnosis literature and their modest computational footprint: a Random
Forest (RF) with $300$ trees; a $k$-nearest-neighbour classifier (KNN, $k=5$);
and a support vector machine (SVM) with a radial-basis-function kernel. To keep
the SVM tractable on the full window set, the kernel was approximated by a
Nystr\"om feature map ($400$ components) followed by a linear SVM. For the
edge-footprint analysis of Section~\ref{sec:results}, two additional
lightweight models---a depth-limited decision tree and a multinomial logistic
regression---were included. All models map a feature descriptor to one of the six
health states.

\subsection{Evaluation Protocols}
\label{subsec:protocols}
The central methodological question is how the labelled windows are partitioned
into training and test sets. Four protocols were defined and applied to identical
features so that differences in reported accuracy are attributable solely to the
partitioning strategy.

\subsubsection{RANDOM}
A stratified random split over all windows, as commonly used in the literature
and by the source baseline. This protocol is retained as a reference precisely
because it is susceptible to leakage: adjacent windows share $50\%$ of their
samples, and all windows of a recording share the same operating point and sensor
coupling, so a model may recognise the originating recording rather than the
fault.

\subsubsection{GROUP}
A leakage-free protocol in which entire recordings are held out. Stratified
group $k$-fold cross-validation ($k=6$) was used, with the recording identifier as
the grouping variable, so that no window from a test recording appears in
training. With six recordings per class, each fold holds out one recording per
class.

\subsubsection{CROSS-SPEED}
A leave-one-supply-frequency-out protocol: the model is trained on two supply
frequencies and tested on the third, and results are averaged over the three
folds. This quantifies generalisation to an unseen operating speed.

\subsubsection{CROSS-LOAD}
A leave-one-load-out protocol: the model is trained on one load condition
(loaded or unloaded) and tested on the other. This quantifies generalisation to
an unseen load.

The GROUP, CROSS-SPEED, and CROSS-LOAD partitions are all group-disjoint at the
recording level and therefore free of the window-overlap leakage that affects the
RANDOM protocol; the latter two additionally introduce a deliberate operating
-condition shift between training and test.

\subsection{Statistical Evaluation}
Classification performance was reported as accuracy and macro-averaged
$F_1$-score, the latter being appropriate for the balanced six-class problem while
remaining sensitive to per-class failures. To account for the limited number of
recordings, the group-aware protocols were repeated over five random seeds, which
reshuffle the fold composition, and the mean and standard deviation over
seeds and folds are reported. Confusion matrices were aggregated from the pooled
out-of-fold predictions of the GROUP protocol. All experiments used fixed seeds
and are reproducible from the released code.

% ---------------------------------------------------------------------------
\begin{table}[t]
\centering
\caption{Per-channel time-domain features. For a window
$x=\{x_i\}_{i=1}^{W}$ with mean $\mu$ and standard deviation $\sigma$.
Features marked $\dagger$ form the six-feature ablation subset.}
\label{tab:features}
\renewcommand{\arraystretch}{1.45}
\begin{tabular}{@{}ll@{}}
\toprule
\textbf{Feature} & \textbf{Definition} \\
\midrule
Mean & $\mu=\frac{1}{W}\sum_i x_i$ \\
Std.\ dev.$^\dagger$ & $\sigma=\sqrt{\frac{1}{W}\sum_i (x_i-\mu)^2}$ \\
Variance & $\sigma^2$ \\
RMS$^\dagger$ & $x_{\mathrm{rms}}=\sqrt{\frac{1}{W}\sum_i x_i^2}$ \\
Peak-to-peak$^\dagger$ & $\max_i x_i-\min_i x_i$ \\
Max.\ absolute & $x_{\max}=\max_i |x_i|$ \\
Skewness$^\dagger$ & $\frac{1}{W}\sum_i (x_i-\mu)^3/\sigma^3$ \\
Kurtosis$^\dagger$ & $\frac{1}{W}\sum_i (x_i-\mu)^4/\sigma^4-3$ \\
Crest factor$^\dagger$ & $x_{\max}/x_{\mathrm{rms}}$ \\
Shape factor & $x_{\mathrm{rms}}/\big(\frac{1}{W}\sum_i |x_i|\big)$ \\
Impulse factor & $x_{\max}/\big(\frac{1}{W}\sum_i |x_i|\big)$ \\
Clearance factor & $x_{\max}/\big(\frac{1}{W}\sum_i \sqrt{|x_i|}\big)^2$ \\
Zero-crossing rate & rate of sign changes of $x_i-\mu$ \\
Amplitude entropy & $-\sum_k p_k\log_2 p_k$ (32-bin histogram) \\
\bottomrule
\end{tabular}
\end{table}
